The hiding cost has two components. The term $c_0\geq 0$ applies to all non-disclosers, regardless of quality. In the French institutional context it captures the regulatory burden imposed by loi ALUR, from the legal exposure if a buyer later discovers a concealed audit via the public ADEME registry or a notaire liability check, to the informational ambiguity that prolongs negotiation and raises the risk that a sale falls through. The term $c_1\geq 0$ is an additional opportunity cost borne only by high-quality sellers ($\theta=1$) who choose to hide, since by not disclosing they forgo the revenue premium that verified high quality would command from EPC-sensitive buyers. Because $c_1>0$, hiding is strictly cheaper for low-quality sellers than for high-quality ones, so the net gain from concealment is larger for $\theta=0$ sellers at any given pool belief, and this differential is the sorting pressure that drives full separation of high-quality sellers in equilibrium (Corollary~\ref{cor:hq_disclose}). Hiding costs are paid unconditionally on sale, reflecting compliance burdens and reputational exposure that arise at listing time rather than at closing.
